\subsection*{Integer Prepairs}

\begin{definitionbox}{Definition}
\begin{definition}[Integer Prepair]
\label{def:integer-prepair}
An \emph{integer prepair} is an ordered pair $(a,b)\in\mathbb{W}\times\mathbb{W}$,
read as the formal difference $a-b$.
\end{definition}
\end{definitionbox}

\begin{remark*}[Standard quantified statement]
\[
(a,b)\in\mathbb{W}\times\mathbb{W}.
\]
\end{remark*}

\begin{remark*}[Predicate reading]
\[
\operatorname{IntegerPrepair}((a,b)).
\]
\end{remark*}

\begin{remark*}[Interpretation]
An integer prepair is the raw ordered-pair datum used to model a formal
difference before quotienting.
\end{remark*}

\begin{dependencies}
\begin{itemize}
  \item \hyperref[def:ordered-pair]{Ordered Pair}
  \item \hyperref[def:whole-numbers]{Whole Numbers}
\end{itemize}
\end{dependencies}

\begin{definitionbox}{Definition}
\begin{definition}[Integer Prepair Set]
\label{def:integer-prepair-set}
The integer prepair set is
\[
\operatorname{Pre}\mathbb{Z}:=\mathbb{W}\times\mathbb{W}.
\]
\end{definition}
\end{definitionbox}

\begin{remark*}[Standard quantified statement]
\[
\operatorname{Pre}\mathbb{Z}=\mathbb{W}\times\mathbb{W}.
\]
\end{remark*}

\begin{remark*}[Predicate reading]
\[
\operatorname{PreIntegerSet}(\operatorname{Pre}\mathbb{Z}).
\]
\end{remark*}

\begin{remark*}[Interpretation]
The preinteger set is the carrier before imposing formal-difference
equivalence.
\end{remark*}

\begin{dependencies}
\begin{itemize}
  \item \hyperref[def:cartesian-product-rel]{Cartesian Product}
  \item \hyperref[def:whole-numbers]{Whole Numbers}
\end{itemize}
\end{dependencies}

\subsection*{Formal-Difference Equivalence}

\begin{definitionbox}{Definition}
\begin{definition}[Formal-Difference Equivalence]
\label{def:formal-difference-equivalence}
For $(a,b),(c,d)\in\operatorname{Pre}\mathbb{Z}$, define
\[
(a,b)\sim_{\mathbb{Z}}(c,d)
\quad:\Longleftrightarrow\quad
a+d=c+b.
\]
\end{definition}
\end{definitionbox}

\begin{remark*}[Standard quantified statement]
\[
\forall (a,b),(c,d)\in\operatorname{Pre}\mathbb{Z},\quad
(a,b)\sim_{\mathbb{Z}}(c,d)
\Longleftrightarrow
a+d=c+b.
\]
\end{remark*}

\begin{remark*}[Predicate reading]
\[
\operatorname{FormalDifferenceEquivalent}((a,b),(c,d)).
\]
\end{remark*}

\begin{remark*}[Interpretation]
Two prepairs represent the same formal difference exactly when cross-addition
balances them.
\end{remark*}

\begin{dependencies}
\begin{itemize}
  \item \hyperref[def:relation]{Relation}
  \item \hyperref[def:integer-prepair-set]{Integer Prepair Set}
  \item \hyperref[def:addition-on-w]{Addition on $\mathbb{W}$}
\end{itemize}
\end{dependencies}

\begin{lemmabox}{Lemma}
\begin{lemma}[Formal-Difference Equivalence Is Reflexive]
\label{lem:formal-difference-equivalence-reflexive}
For every $(a,b)\in\operatorname{Pre}\mathbb{Z}$,
\[
(a,b)\sim_{\mathbb{Z}}(a,b).
\]
\hyperref[prf:formal-difference-equivalence-reflexive]{\textit{Go to proof.}}
\end{lemma}
\end{lemmabox}

\begin{remark*}[Standard quantified statement]
\[
\forall (a,b)\in\operatorname{Pre}\mathbb{Z},\ (a,b)\sim_{\mathbb{Z}}(a,b).
\]
\end{remark*}

\begin{remark*}[Predicate reading]
\[
\operatorname{Reflexive}(\sim_{\mathbb{Z}},\operatorname{Pre}\mathbb{Z}).
\]
\end{remark*}

\begin{remark*}[Interpretation]
Every integer prepair represents the same formal difference as itself.
\end{remark*}

\begin{dependencies}
\begin{itemize}
  \item \hyperref[def:formal-difference-equivalence]{Formal-Difference Equivalence}
  \item \hyperref[def:addition-on-w]{Addition on $\mathbb{W}$}
\end{itemize}
\end{dependencies}

\begin{lemmabox}{Lemma}
\begin{lemma}[Formal-Difference Equivalence Is Symmetric]
\label{lem:formal-difference-equivalence-symmetric}
For all prepairs $(a,b),(c,d)$,
\[
(a,b)\sim_{\mathbb{Z}}(c,d)
\Longrightarrow
(c,d)\sim_{\mathbb{Z}}(a,b).
\]
\hyperref[prf:formal-difference-equivalence-symmetric]{\textit{Go to proof.}}
\end{lemma}
\end{lemmabox}

\begin{remark*}[Standard quantified statement]
\[
\forall p,q\in\operatorname{Pre}\mathbb{Z},\quad
p\sim_{\mathbb{Z}}q \Longrightarrow q\sim_{\mathbb{Z}}p.
\]
\end{remark*}

\begin{remark*}[Predicate reading]
\[
\operatorname{Symmetric}(\sim_{\mathbb{Z}},\operatorname{Pre}\mathbb{Z}).
\]
\end{remark*}

\begin{remark*}[Interpretation]
Formal-difference equivalence is reversible.
\end{remark*}

\begin{dependencies}
\begin{itemize}
  \item \hyperref[def:formal-difference-equivalence]{Formal-Difference Equivalence}
  \item \hyperref[def:addition-on-w]{Addition on $\mathbb{W}$}
\end{itemize}
\end{dependencies}

\begin{lemmabox}{Lemma}
\begin{lemma}[Formal-Difference Equivalence Is Transitive]
\label{lem:formal-difference-equivalence-transitive}
For all prepairs $(a,b),(c,d),(e,f)$,
\[
(a,b)\sim_{\mathbb{Z}}(c,d)\land
(c,d)\sim_{\mathbb{Z}}(e,f)
\Longrightarrow
(a,b)\sim_{\mathbb{Z}}(e,f).
\]
\hyperref[prf:formal-difference-equivalence-transitive]{\textit{Go to proof.}}
\end{lemma}
\end{lemmabox}

\begin{remark*}[Standard quantified statement]
\[
\forall p,q,r\in\operatorname{Pre}\mathbb{Z},\quad
p\sim_{\mathbb{Z}}q\land q\sim_{\mathbb{Z}}r
\Longrightarrow
p\sim_{\mathbb{Z}}r.
\]
\end{remark*}

\begin{remark*}[Predicate reading]
\[
\operatorname{Transitive}(\sim_{\mathbb{Z}},\operatorname{Pre}\mathbb{Z}).
\]
\end{remark*}

\begin{remark*}[Interpretation]
Formal-difference equivalence can be chained through an intermediate prepair.
\end{remark*}

\begin{dependencies}
\begin{itemize}
  \item \hyperref[thm:addition-on-w-cancellation]{Addition on $\mathbb{W}$ Satisfies Cancellation}
\end{itemize}
\end{dependencies}

\begin{theorembox}{Theorem}
\begin{theorem}[Formal-Difference Equivalence Is an Equivalence Relation]
\label{thm:formal-difference-equivalence-relation}
The relation $\sim_{\mathbb{Z}}$ is an equivalence relation on
$\operatorname{Pre}\mathbb{Z}$.
\hyperref[prf:formal-difference-equivalence-relation]{\textit{Go to proof.}}
\end{theorem}
\end{theorembox}

\begin{remark*}[Standard quantified statement]
\[
\operatorname{EquivalenceRelation}(\sim_{\mathbb{Z}},\operatorname{Pre}\mathbb{Z}).
\]
\end{remark*}

\begin{remark*}[Predicate reading]
\[
\operatorname{EquivalenceRelation}(\sim_{\mathbb{Z}},\operatorname{Pre}\mathbb{Z}).
\]
\end{remark*}

\begin{remark*}[Interpretation]
The reflexive, symmetric, and transitive laws make quotienting by
$\sim_{\mathbb{Z}}$ legitimate.
\end{remark*}

\begin{dependencies}
\begin{itemize}
\item 
\hyperref[def:formal-difference-equivalence]{Formal-Difference Equivalence}
  \item \hyperref[lem:formal-difference-equivalence-reflexive]{Formal-Difference Equivalence Is Reflexive}
  \item \hyperref[lem:formal-difference-equivalence-symmetric]{Formal-Difference Equivalence Is Symmetric}
  \item \hyperref[lem:formal-difference-equivalence-transitive]{Formal-Difference Equivalence Is Transitive}
\end{itemize}
\end{dependencies}

\subsection*{\texorpdfstring{Definition of $\mathbb{Z}$}{Definition of Z}}

\begin{definitionbox}{Definition}
\begin{definition}[Integers]
\label{def:integers}
The integers are the quotient
\[
\mathbb{Z}:=\operatorname{Pre}\mathbb{Z}/{\sim_{\mathbb{Z}}}.
\]
\end{definition}
\end{definitionbox}

\begin{remark*}[Standard quantified statement]
\[
\mathbb{Z}=\operatorname{Pre}\mathbb{Z}/{\sim_{\mathbb{Z}}}.
\]
\end{remark*}

\begin{remark*}[Predicate reading]
\[
\mathbb{Z}=\operatorname{Quotient}(\operatorname{Pre}\mathbb{Z},\sim_{\mathbb{Z}}).
\]
\end{remark*}

\begin{remark*}[Interpretation]
The integers are constructed as a quotient of formal differences by
formal-difference equivalence.
\end{remark*}

\begin{dependencies}
\begin{itemize}
  \item \hyperref[def:quotient-set]{Quotient Set}
  \item \hyperref[def:integer-prepair-set]{Integer Prepair Set}
  \item \hyperref[def:formal-difference-equivalence]{Formal-Difference Equivalence}
\end{itemize}
\end{dependencies}

\begin{definitionbox}{Definition}
\begin{definition}[Integer Class Notation]
\label{def:integer-class-notation}
The notation $[a,b]$ denotes the $\sim_{\mathbb{Z}}$-equivalence class of the
prepair $(a,b)$.
\end{definition}
\end{definitionbox}

\begin{remark*}[Standard quantified statement]
\[
[a,b]=[(a,b)]_{\sim_{\mathbb{Z}}}.
\]
\end{remark*}

\begin{remark*}[Predicate reading]
\[
\operatorname{EquivalenceClassNotation}([a,b],(a,b),\sim_{\mathbb{Z}}).
\]
\end{remark*}

\begin{remark*}[Interpretation]
The bracket notation names the equivalence class of an integer prepair.
\end{remark*}

\begin{dependencies}
\begin{itemize}
  \item \hyperref[def:integers]{Integers}
  \item \hyperref[def:formal-difference-equivalence]{Formal-Difference Equivalence}
\end{itemize}
\end{dependencies}

\subsection*{\texorpdfstring{Zero and One in $\mathbb{Z}$}{Zero and One in Z}}

\begin{definitionbox}{Definition}
\begin{definition}[Zero in $\mathbb{Z}$]
\label{def:zero-in-z}
Define
\[
0_{\mathbb{Z}}:=[0,0].
\]
\end{definition}
\end{definitionbox}

\begin{remark*}[Standard quantified statement]
\[
0_{\mathbb{Z}}=[0,0].
\]
\end{remark*}

\begin{remark*}[Predicate reading]
\[
\operatorname{ZeroElement}(0_{\mathbb{Z}},\mathbb{Z}).
\]
\end{remark*}

\begin{remark*}[Interpretation]
The integer zero is the class of the zero formal difference.
\end{remark*}

\begin{dependencies}
\begin{itemize}
  \item \hyperref[def:integer-class-notation]{Integer Class Notation}
  \item \hyperref[def:whole-numbers]{Whole Numbers}
\end{itemize}
\end{dependencies}

\begin{definitionbox}{Definition}
\begin{definition}[One in $\mathbb{Z}$]
\label{def:one-in-z}
Define
\[
1_{\mathbb{Z}}:=[1,0].
\]
\end{definition}
\end{definitionbox}

\begin{remark*}[Standard quantified statement]
\[
1_{\mathbb{Z}}=[1,0].
\]
\end{remark*}

\begin{remark*}[Predicate reading]
\[
\operatorname{OneElement}(1_{\mathbb{Z}},\mathbb{Z}).
\]
\end{remark*}

\begin{remark*}[Interpretation]
The integer one is the class of the whole-number prepair representing one.
\end{remark*}

\begin{dependencies}
\begin{itemize}
  \item \hyperref[def:integer-class-notation]{Integer Class Notation}
  \item \hyperref[def:whole-numbers]{Whole Numbers}
\end{itemize}
\end{dependencies}

\begin{theorembox}{Theorem}
\begin{theorem}[Zero Is an Integer]
\label{thm:zero-is-an-integer}
The class $0_{\mathbb{Z}}$ is an element of $\mathbb{Z}$.
\hyperref[prf:zero-is-an-integer]{\textit{Go to proof.}}
\end{theorem}
\end{theorembox}

\begin{remark*}[Standard quantified statement]
\[
0_{\mathbb{Z}}\in\mathbb{Z}.
\]
\end{remark*}

\begin{remark*}[Predicate reading]
\[
\operatorname{ElementOf}(0_{\mathbb{Z}},\mathbb{Z}).
\]
\end{remark*}

\begin{remark*}[Interpretation]
The distinguished zero class belongs to the quotient carrier.
\end{remark*}

\begin{dependencies}
\begin{itemize}
  \item \hyperref[def:integers]{Integers}
  \item \hyperref[def:zero-in-z]{Zero in $\mathbb{Z}$}
  \item \hyperref[def:integer-class-notation]{Integer Class Notation}
\end{itemize}
\end{dependencies}

\begin{theorembox}{Theorem}
\begin{theorem}[One Is an Integer]
\label{thm:one-is-an-integer}
The class $1_{\mathbb{Z}}$ is an element of $\mathbb{Z}$.
\hyperref[prf:one-is-an-integer]{\textit{Go to proof.}}
\end{theorem}
\end{theorembox}

\begin{remark*}[Standard quantified statement]
\[
1_{\mathbb{Z}}\in\mathbb{Z}.
\]
\end{remark*}

\begin{remark*}[Predicate reading]
\[
\operatorname{ElementOf}(1_{\mathbb{Z}},\mathbb{Z}).
\]
\end{remark*}

\begin{remark*}[Interpretation]
The distinguished one class belongs to the quotient carrier.
\end{remark*}

\begin{dependencies}
\begin{itemize}
  \item \hyperref[def:integers]{Integers
}
  \item \hyperref[def:one-in-z]{One in $\mathbb{Z}$}
  \item \hyperref[def:integer-class-notation]{Integer Class Notation}
\end{itemize}
\end{dependencies}

\begin{theorembox}{Theorem}
\begin{theorem}[Zero Is Not One in $\mathbb{Z}$]
\label{thm:zero-not-one-in-z}
In $\mathbb{Z}$,
\[
0_{\mathbb{Z}}\ne 1_{\mathbb{Z}}.
\]
\hyperref[prf:zero-not-one-in-z]{\textit{Go to proof.}}
\end{theorem}
\end{theorembox}

\begin{remark*}[Standard quantified statement]
\[
0_{\mathbb{Z}}\ne 1_{\mathbb{Z}}.
\]
\end{remark*}

\begin{remark*}[Predicate reading]
\[
\operatorname{Distinct}(0_{\mathbb{Z}},1_{\mathbb{Z}}).
\]
\end{remark*}

\begin{remark*}[Interpretation]
The quotient construction preserves the distinction between the zero and one
integer classes.
\end{remark*}

\begin{dependencies}
\begin{itemize}
  \item \hyperref[thm:zero-not-in-n]{$0\notin\mathbb{N}$}
\end{itemize}
\end{dependencies}
